\scanpage{35}\mspage{34}
Interprétons la donnée d'une $F'\preccurlyeq F$ induisant $G'\preccurlyeq G$, en termes des données initiales, exprimées par un diagramme d'applications (où $S=|F|$, $T=|G|$, $I$ et $J$ les ens \inserted{ord} d'indices pour $F,G$, $J'$ celui pour $G'$)

\diagram{\begin{tikzcd}[ampersand replacement=\&,column sep=huge,row sep=large]
\& \deleted{\varphi}f^{-1}(J)=T \arrow[r,hook,no head]\arrow[d,"{\deleted{\unsure{\psi|\varphi}}\unsure{g_J^{*}|g_J}}"]\arrow[dd,bend right=30,"{\deleted{\varphi}\unsure{f_J^{*}|f_J}=\deleted{\varphi}f\mid T}"'] \& S\arrow[dd,"{ \deleted{\varphi}f\ \text{surj.}}"]\\
(*) \& J'\arrow[d,"{\deleted{\chi}h_J}"] \& {}\\
\& J\arrow[r,hook,no head,"i","\substack{\text{partie}\\\text{fermée de}\\I}"'] \& I
\end{tikzcd}}
\margin{\uline{NB} On part de $f\colon S\mathbin{-}I$ et \deleted{factorise} \inserted{$J\subset I$}, d'où $T=f^{-1}(J)$, $f_J\colon T\longrightarrow J$, et on factorise $f_J$ en $T\mathbin{-}J'\mathbin{-}J$}

La donnée de $G'$ équivaut à celle d'une factorisation de $\varphi'=\varphi\mid T$ en $T\xrightarrow{\deleted{\unsure{\psi'}}\unsure{g_J|g_J^{*}}}J'\xrightarrow{\deleted{\unsure{\chi'}}\unsure{h_J^{*}|h_J}}J$, où $\chi'$ est une \deleted{\unsure{un}} appl. croiss. surjective d'ens ordonnés, et $\psi'$ une application ensembliste surjection. La donnée d'une $F'$ équivaut à une factorisation analogue pour $\varphi$ en
\diagram{\begin{tikzcd}[ampersand replacement=\&,column sep=large]
S\arrow[r,"{\deleted{\unsure{\psi}}g}"] \& I'\arrow[r,"{\deleted{\chi}h}"] \& I \& \substack{g\ \text{surj}\\ h\ \text{croiss. surj.}}
\end{tikzcd}}

Dire que ça induit \deleted{\unsure{un diag}} la \deleted{\unsure{situation}} \inserted{sous-d. donnée $G'$ de $G$} donnée signifie qu'\deleted{\unsure{elle}} \deleted{\unsure{a des d'applications}} \inserted{existe $i'\colon J'\hookrightarrow I'$} tel que les deux carrés commutent
\diagram{\begin{tikzcd}[ampersand replacement=\&,column sep=huge]
T\arrow[r,hook,no head]\arrow[d,"g_J"]\arrow[dd,bend right=30,"f_J"'] \& S\arrow[d,"g"]\\
J'\arrow[r,hook,"i'"]\arrow[d,"h_J"] \& I'\arrow[d,"h"] \& \text{avec }\deleted{\unsure{\xi}}J'=\deleted{\unsure{\chi|k}}h^{-1}(J)\\
J\arrow[r,hook,"i"] \& I
\end{tikzcd}}
